\clearpage
\VerdierPhysicalPageAnchor{0187}%
\begingroup\small\setlength{\parskip}{0pt}\setlength{\emergencystretch}{1.5em}
\setlength{\abovedisplayskip}{4pt plus 2pt minus 2pt}%
\setlength{\belowdisplayskip}{4pt plus 2pt minus 2pt}%

\VerdierSemanticAnchor{ch3d-2-1-1-base-change}
现在设 $m\colon Z\longrightarrow X$ 是复形的一个态射。复形 $Z'$ 显然是
$Z$ 上由沿 $Z\xrightarrow{m}X$ 的基变换得到的对象，并由下图中左侧的竖直箭头给出：

\VerdierSemanticAnchor{ch3d-2-1-1-base-change-square}
\begin{center}
\begin{tikzcd}[column sep=4.2em,row sep=2.6em]
Z' \arrow[r] \arrow[d] & X(u,i) \arrow[d,"{p(u,i)}"] \\
Z \arrow[r,"m"'] & X
\end{tikzcd}
\end{center}
该箭头同构于一个形如 $p(w,i)\colon Z(w,i)\longrightarrow Z$ 的复形态射，其中
$w$ 是适当的态射。

\VerdierSemanticAnchor{ch3d-2-1-2}
\medskip\noindent\textbf{引理 2.1.2.}\quad
\textit{a) 态射
$H^j(p(u,i))\colon H^j(X(u,i))\longrightarrow H^j(X)$ 在
$j\neq i,i+1$ 时是同构。}

\noindent\textit{b) 有正合列：}

\VerdierSemanticAnchor{ch3d-2-1-2-exact-sequences}
\[
0\longrightarrow H^i(X(u,i))\longrightarrow H^i(X)
\longrightarrow
\frac{\operatorname{Ker}(d_X^i)+\operatorname{Im}(u)}
     {\operatorname{Im}(d_X^{i-1})+\operatorname{Im}(u)}
\longrightarrow0,
\]
\[
0\longrightarrow
\frac{\operatorname{Ker}(d_X^i)+\operatorname{Im}(u)}
     {\operatorname{Im}(d_X^{i-1})+\operatorname{Im}(u)}
\longrightarrow
\frac{X^i}{\operatorname{Im}(d_X^{i-1})+\operatorname{Im}(u)}
\longrightarrow
\frac{X^i}{\operatorname{Ker}(d_X^i)+\operatorname{Im}(u)}
\longrightarrow0,
\]
\[
0\longrightarrow
\frac{X^i}{\operatorname{Ker}(d_X^i)+\operatorname{Im}(u)}
\longrightarrow H^{i+1}(X(u,i))\longrightarrow H^{i+1}(X)
\longrightarrow0.
\]

\noindent\textit{c) $p(u,i)$ 是拟同构的充分必要条件为}
\[
\operatorname{Im}(d_X^{i-1})+\operatorname{Im}(u)=X^i.
\]
\VerdierSemanticAnchor{ch3d-2-1-2-proof}
图追逐留给读者。

\VerdierSemanticAnchor{ch3d-2-1-3}
\medskip\noindent\textbf{定义 2.1.3.}\quad
复形态射 $m\colon Z\longrightarrow X$ 称为一个\emph{高度为 $i$ 的初等态射}，
如果它在 $\mathcal{F}\ell_1(\operatorname{comp}(\mathcal A))$ 中同构于形如
$p(u,i)$ 的态射，其中 $u$ 是满态射。

\VerdierSemanticAnchor{ch3d-2-1-3-properties}
因此初等态射特别是拟同构（(2.1.2), c)）。初等态射是
$\operatorname{comp}(\mathcal A)$ 中的满态射。
$\operatorname{comp}(\mathcal A)$ 中的满态射是初等态射，当且仅当其核是
至多含有两个非零分量的无环复形。同一高度的两个初等态射的复合是
\endgroup
